\begin{table*}[t]
\centering
\caption{Real-data learner and strategy decomposition. Entries are equal-dataset mean paired differences in percentage points; the 20 paired split differences are first averaged within each dataset and condition. These are descriptive pipeline contrasts, not population inference.}
\label{tab:real-decomposition-summary}
\small
\setlength{\tabcolsep}{3.5pt}
\renewcommand{\arraystretch}{1.08}
\begin{tabularx}{\textwidth}{llr*{3}{>{\centering\arraybackslash}X}}
\toprule
Metric & Design & Budget & Learner gain & CS addition & RN addition \\
\midrule
AUROC & OS & 0.05 & +2.540 & -0.701 & +0.144 \\
AUROC & OS & 0.20 & +2.043 & -0.159 & -0.638 \\
AUROC & TS & 0.05 & +2.585 & -0.448 & +0.118 \\
AUROC & TS & 0.20 & +2.042 & -0.153 & -0.592 \\
\midrule
AP & OS & 0.05 & +2.188 & -1.113 & +1.207 \\
AP & OS & 0.20 & +2.256 & -0.285 & -0.518 \\
AP & TS & 0.05 & +2.240 & -0.707 & +1.094 \\
AP & TS & 0.20 & +2.064 & -0.204 & -0.259 \\
\bottomrule
\end{tabularx}
\par\medskip
\begin{minipage}{\textwidth}\small
Learner gain: naive CatBoost minus logistic probability correction. CS/RN additions: the respective CatBoost strategy minus naive CatBoost. Positive values favour the first named method. Strategy-specific tuning is retained; no test-selected representative or group maximum is used.
\end{minipage}
\end{table*}
